% ============================================================
%  第三章  因果图与因果效应的识别
% ============================================================
\chapter{因果图与因果效应的识别}

\begin{flushright}
\kaishu\small
眼睛精确地遵从心的行动。\\[2pt]
——爱默生（Emerson，1860）
\end{flushright}

\section*{序言}
\addcontentsline{toc}{section}{序言}

上一章讨论的是从原始数据中学习因果关系的方法。本章探讨如何从数据与领域中被认为合理的定性因果假设相结合，来推断这类关系。更广泛地说，本章旨在帮助研究者交流关于因果关系的定性假设、阐明这些假设的推演后果，并从假设、实验与数据的结合中得出因果推断。我们的主要任务是判定：给定的假设是否足以从非实验数据中评估因果效应的强度。

因果效应使我们能够预测系统将如何响应假设性的干预——例如政策决定或日常活动中的行动。如第一章（\S 1.3）所见，这类预测是经验科学的标志，且无法仅从概率信息中识别出来；它们立足于因果关系，事实上也定义了因果关系。本章用因果图给“干预”概念赋予形式语义，并给出用干预前概率表述干预后概率的显式公式。其含义是：只要数据补充以一个既无环、又不含潜变量的因果图，每一种干预的效果都可以由非实验数据估计。若某些变量未被测量，便产生了可识别性问题；本章为此发展了一个非参数框架，用以分析一般因果关系的识别、尤其是因果效应的识别。我们将看到，因果图在这一分析中提供了强大的数学工具；只需极其简单的检验，就可以查询这些图，判定现有假设是否足以识别因果效应。若足以识别，图便给出用被观测分布表述因果效应的数学表达式；若不足以识别，则可通过查询图来提示还需要哪些额外观测或辅助实验，方能获得所需的推断。

因果效应的图形分析还引出另一件工具——\emph{干预演算}：一组推理规则，用以把涉及干预与观测的语句变换为其他此类语句，从而提供一种句法方法，来导出（或验证）关于干预及其与观测如何相互作用的断言。借助这一演算，读者将能够：(i) 在数学上判定给定的一组协变量是否适于控制混杂；(ii) 处理位于因果通路上的测量；(iii) 用一组测量替换另一组。

最后，我们将说明这一新演算如何澄清那些曾在哲学家、统计学家、经济学家与心理学家之间引发争议与误解的概念，包括结构方程与回归方程之别、直接效应与间接效应的定义，以及结构方程与 Neyman–Rubin 模型之间的关系。

\section{引论}

本章所处理的问题，最好通过 Cochran 的一个经典例子来说明（见 Wainer 1989）。考虑一个用土壤熏蒸剂（$X$）通过控制线虫种群（$Z$）来提高燕麦产量（$Y$）的实验；除控制线虫之外，熏蒸剂还可能对产量有直接效应（有益的和不利的皆有）。当这一典型研究被若干因素复杂化时，我们希望评估熏蒸剂对产量的总效应。第一，受控随机实验不可行——农夫坚持自己决定哪些地块要熏蒸。第二，农夫对处理的选择取决于去年的线虫种群（$Z_0$），而这是一个与今年种群强相关的未知量。于是我们遇到一个经典的混杂偏倚情形，它妨碍对处理效应的评估，且不论样本量多大都无济于事。所幸，通过对土壤样本的实验室分析，我们可以确定处理前后的线虫种群；此外，由于已知熏蒸剂只在短期内有效，我们可以稳妥地假定它们不影响处理后存活线虫的生长。线虫的生长反而取决于鸟类（及其他捕食者）的种群，而后者与去年的线虫种群相关，因而与处理本身相关。

本章所发展的方法，使研究者能够把这类复杂的考量翻译为一种形式语言，从而促进下列任务：
\begin{enumerate}[label=\arabic*.,leftmargin=1.8em]
\item 阐明模型背后的假设；
\item 判定这些假设是否足以获得目标量（熏蒸剂对产量的总效应）的一致估计；
\item 若第 2 项答案为肯定，则给出目标量用被观测量的分布表示的闭式表达式；
\item 若第 2 项答案为否定，则提示一组观测与实验，若付诸实施便可使一致估计成为可能。
\end{enumerate}

分析的第一步是构造一个因果图，如图 3.1 所示，它代表研究者对本领域中可测量量之间主要因果影响的理解。例如，$Z_1,Z_2,Z_3$ 分别表示处理前、处理后与季末的线虫种群。$Z_0$ 表示去年的线虫种群；由于它是未知量，故以空心圆表示，量 $B$（鸟类及其他捕食者的种群）亦然。图中的连线有两种：发自未测量量的连线用虚线箭头表示，连接被测量量的连线用实线箭头表示。图中所体现的实质性假设是一些否定的因果断言，它们通过图中\emph{缺失}的连线传达。例如，$Z_1$ 与 $Y$ 之间缺失箭头，表示研究者认为处理前的线虫不可能直接影响燕麦植株；它们对燕麦产量的全部影响都以处理后状况 $Z_2$、$Z_3$ 为中介。我们的目的不是证实或否证这类领域特定的假设，而是检验给定的一组假设是否足以从非实验数据中量化因果效应——在此即估计熏蒸剂对产量的总效应。

图 3.1 的因果图在许多方面与 Wright (1921) 所创的路径图相似：二者都反映研究者对本领域因果影响的主观定性知识，都采用有向无环图，都允许纳入潜变量或未测量量。主要差别在于分析方法。

\begin{figure}[htbp]
\centering
\begin{tikzpicture}[causal,scale=1.0,every node/.style={transform shape}]
  \node[vtx] (Z0) at (0,3.0) {$Z_0$};
  \node[vtx] (Z1) at (-1.35,1.85) {$Z_1$};
  \node[vtx] (B)  at (1.55,2.05) {$B$};
  \node[vtx] (X)  at (-2.35,0.95) {$X$};
  \node[vtx] (Z2) at (0.0,0.95) {$Z_2$};
  \node[vtx] (Z3) at (1.95,0.95) {$Z_3$};
  \node[vtx] (Y)  at (0.35,-0.35) {$Y$};
  \draw[edd] (Z0)--(Z1); \draw[edd] (Z0)--(B); \draw[edd] (Z0)--(X);
  \draw[edd] (B)--(Z3);
  \draw[ed] (Z1)--(Z2); \draw[ed] (X)--(Z2);
  \draw[ed] (Z2)--(Z3);
  \draw[ed] (X)--(Y); \draw[ed] (Z2)--(Y); \draw[ed] (Z3)--(Y);
\end{tikzpicture}
\caption{表示熏蒸剂（$X$）对产量（$Y$）之效应的因果图。}
\label{fig:3.1}
\end{figure}

第一，路径图大多是在带高斯噪声的线性模型语境中分析的，而因果图允许任意的非线性交互。事实上，我们对因果效应的分析将完全是非参数的，不对方程与分布的函数形式作任何承诺。第二，因果图不仅被用作交流假设的被动语言，还被用作导出所需量的主动计算装置。例如，将要描述的方法使研究者能够审视图 3.1，并立刻得出以下结论：
\begin{enumerate}[label=\arabic*.,leftmargin=1.8em]
\item $X$ 对 $Y$ 的总效应可以由 $X,Z_1,Z_2,Z_3,Y$ 的被观测分布一致地估计。
\item $X$ 对 $Y$ 的总效应（假定一切变量均为离散）由下述公式给出\footnote{第一章使用记号 $P_x(y)$；因下标处理不便，此后改用 $P(y\mid\hat{x})$ 或 $P(y\mid\doxp(x))$。读者若觉得 \textup{(3.1)} 陌生，不必畏惧；读完 \S 3.4 之后，读者应能比解代数方程更轻松地导出这类公式。注意 $x'$ 只是求和指标，取遍 $X$ 的各个值。}：
\begin{equation}
\begin{aligned}
P(y\mid\hat{x}) = {}& \sum_{z_1}\sum_{z_2}\sum_{z_3} P(y\mid z_2,z_3,x)\,P(z_2\mid z_1,x)\\
& \times \sum_{x'} P(z_3\mid z_1,z_2,x')\,P(z_1,x'),
\end{aligned} \tag{3.1}
\end{equation}
其中 $P(y\mid\hat{x})$ 表示在外来干预把处理设定为水平 $X=x$ 的条件下，产量达到 $Y=y$ 的概率。
\item 若 $Y$ 与 $Z_3$ 混杂，则无法一致地估计 $X$ 对 $Y$ 的总效应；然而 $Z_2$ 与 $Y$ 的混杂不会使 $P(y\mid\hat{x})$ 的公式失效。
\end{enumerate}
这些结论或可通过分析图的图形性质得到，或可通过执行一系列（受图支配的）符号推导得到，而后者正是产生 \textup{(3.1)} 这类因果效应公式的途径。

\section{马尔可夫模型中的干预}

\subsection{图作为干预的模型}

在第一章（\S 1.3）我们看到，因果模型不同于概率模型，可用于预测干预的效果。这一附加功能要求联合分布 $P$ 补充以一个因果图——即一个标明所关心变量之间因果联系的有向无环图 $G$。本节详述干预的性质，并给出其效果的显式公式。

DAG 的因果解读与关联解读之间的联系，通过以机制为基础的因果论而形成，其根源可追溯到计量经济学的早期工作（Frisch 1938; Haavelmo 1943; Simon 1953）。在这一说明中，关于因果影响的断言（如图 3.1 中各连线所指定的那些）代表相应量之间自主的物理机制；这些机制被表示为受随机扰动扰动的函数关系。呼应这一传统，Pearl and Verma (1991) 用函数关系而非概率关系来解释 DAG 的因果解读（见 \textup{(1.40)} 与定义 2.2.2）；换言之，DAG $G$ 中每个子—父家族都代表一个确定性函数
\begin{equation}
x_i = f_i(\pa_i, e_i),\qquad i=1,\dots,n, \tag{3.2}
\end{equation}
其中 $\pa_i$ 是变量 $X_i$ 在 $G$ 中的父节点；$e_i\ (1\le i\le n)$ 是联合独立、分布任意的随机扰动。这些扰动项代表研究者选择不纳入分析的独立背景因素。若其中任一因素被判断为影响两个或更多变量（从而违反独立性假设），则该因素必须作为未测量（或潜）变量进入分析，并在图中以空心节点表示，如图 3.1 中的 $Z_0$ 与 $B$。例如，图 3.1 的模型所传达的因果假设对应于下述方程组：
\begin{equation}
\begin{aligned}
Z_0 &= f_0(e_0), &\quad B &= f_B(Z_0,e_B), &\quad Z_1 &= f_1(Z_0,e_1),\\
X &= f_X(Z_0,e_X), &\quad Z_2 &= f_2(X,Z_1,e_2),\\
Z_3 &= f_3(B,Z_2,e_3), &\quad Y &= f_Y(X,Z_2,Z_3,e_Y).
\end{aligned} \tag{3.3}
\end{equation}

更一般地，我们可以把所有未观测因素（包括 $e_i$）合并为一个背景变量集 $U$，再用分布函数 $P(u)$——或 $P(u)$ 的某些方面（如独立关系）——概括其特征。于是，因果模型的完整指定包含两个组成部分：一组函数关系
\begin{equation}
x_i = f_i(\pa_i, u_i),\qquad i=1,\dots,n, \tag{3.4}
\end{equation}
以及背景因素上的联合分布函数 $P(u)$。若与因果模型 $M$ 相联系的图 $G(M)$ 无环，则称 $M$ 为\emph{半马尔可夫的}。若此外背景变量相互独立，则称 $M$ 为\emph{马尔可夫的}，因为所得的被观测变量分布将相对于 $G(M)$ 为马尔可夫的（见定理 1.4.1）。因此，若被观测变量为 $\{X,Y,Z_1,Z_2,Z_3\}$，则图 3.1 所描述的模型是半马尔可夫的；若 $Z_0$ 与 $B$ 也被观测，它便成为马尔可夫的。第七章将分析一般的非马尔可夫模型，但本章假定一切模型或是马尔可夫的，或是带有未观测变量的马尔可夫模型（即半马尔可夫的）。

无须说，我们很少能掌握 $P(u)$ 乃至 $f_i$。然而，阐明一个完全指定模型的数学内容仍然重要，因为唯有如此，才能从部分指定的模型（如图 3.1 所描述者）中作出有效推断。

方程模型 \textup{(3.2)} 是所谓结构方程模型（Wright 1921; Goldberger 1973）的非参数对应物，只是方程的函数形式（以及扰动项的分布）仍未被指定。结构方程中的等号传达“由……所决定”这一不对称的反事实关系，每个方程代表一个稳定自主的机制。例如，关于 $Y$ 的方程表明：无论我们当前对 $Y$ 观测到什么，也无论其他方程可能发生什么变化，若变量 $(X,Z_2,Z_3,e_Y)$ 分别取值 $(x,z_2,z_3,e_Y)$，则 $Y$ 将取函数 $f_Y$ 所规定的值。

回想 \S 1.4 的讨论，每个子—父关系的函数式刻画导致与贝叶斯网络相同的联合分布递归分解：
\begin{equation}
P(x_1,\dots,x_n) = \prod_i P(x_i\mid\pa_i), \tag{3.5}
\end{equation}
在图 3.1 的例子中，它给出
\begin{equation}
\begin{aligned}
P(z_0,x,z_1,b,z_2,z_3,y) = {}& P(z_0)P(x\mid z_0)P(z_1\mid z_0)P(b\mid z_0)\\
& \times P(z_2\mid x,z_1)P(z_3\mid z_2,b)P(y\mid x,z_2,z_3).
\end{aligned} \tag{3.6}
\end{equation}
此外，函数式刻画提供了一种便利的语言，用以指定所得分布将如何响应外部干预而变化。其做法是把每个干预编码为对一组精选函数的改动，同时保持其他函数不变。一旦知道了干预所改变的机制是哪些、以及改变的性质，干预的总体效果便可通过修改模型中相应的方程、并用修改后的模型计算新的概率函数来预测。

最简单的一类外部干预，是迫使单个变量（比如 $X_i$）取某个定值 $x_i$。这种我们称为“原子”干预的操作，相当于把 $X_i$ 从旧函数机制 $x_i=f_i(\pa_i,u_i)$ 的影响下抬起，置于一个新机制的影响之下：该机制设定取值 $x_i$，而所有其他机制不受扰动。形式上，这一原子干预——记作 $\doxp(X_i=x_i)$，或简记 $\doxp(x_i)$\footnote{一个等价的记号，用集合 $(x)$ 代替 $\doxp(x)$，见于 Pearl (1995a)。$\doxp(x)$ 记号最早用于 Goldszmidt and Pearl (1992)，如今日益流行。Lauritzen (2001) 使用 $P(y\mid X\backslash x)$。表达式 $P(y\mid\doxp(x))$ 在意图上等价于 Neyman (1923) 与 Rubin (1974) 所引入潜在结局模型中的 $P(Y_x=y)$，也等价于 Lewis (1973b) 反事实理论中的 $P[(X=x)\ \square\!\!\to (Y=y)]$。这些概念的语义基础将在 \S 3.6.3 与第七章讨论。}——相当于从模型中删去方程 $x_i=f_i(\pa_i,u_i)$，并在其余方程中代入 $X_i=x_i$。如此产生的新模型表示系统在干预 $\doxp(X_i=x_i)$ 之下的行为；解出 $X_j$ 的分布，便得到 $X_i$ 对 $X_j$ 的因果效应，记作 $P(x_j\mid\hat{x}_i)$。更一般地，当干预迫使变量子集 $X$ 取定值 $x$ 时，就要从 \textup{(3.4)} 的模型中剪除一个方程子集——$X$ 的每个成员对应一个——从而在其余变量上定义一个新的分布，它完全刻画了干预的效果。\footnote{把干预视为方程修改者的基本观点源自 Marschak (1950) 与 Simon (1953)。把干预显式翻译为从模型中“抹去”方程的做法，最早由 Strotz and Wold (1960) 提出，后为 Fisher (1970) 与 Sobel (1990) 所用。这一翻译的图形推演最先由 Spirtes 等 (1993)、继由 Pearl (1993b) 阐明。Balke and Pearl (1994a,b) 把它提出来作为定义反事实的基础；见 \textup{(3.51)}。}

\begin{definition}[因果效应]
给定两个不相交的变量集 $X$ 与 $Y$，$X$ 对 $Y$ 的\emph{因果效应}（记作 $P(y\mid\hat{x})$ 或 $P(y\mid\doxp(x))$）是一个从 $X$ 到 $Y$ 上概率分布空间的函数。对 $X$ 的每个实现 $x$，$P(y\mid\hat{x})$ 给出：从 \textup{(3.4)} 的模型中删去与 $X$ 中变量相对应的全部方程、并在其余方程中代入 $X=x$ 之后，所诱导出的 $Y=y$ 的概率。
\end{definition}

显然，与约简后方程组相对应的图，是 $G$ 的一个删去了全部进入 $X$ 之箭头的子图（Spirtes et al. 1993）。差值 $E(Y\mid\doxp(x'))-E(Y\mid\doxp(x''))$ 有时被取作“因果效应”的定义（Rosenbaum and Rubin 1983），其中 $x'$ 与 $x''$ 是 $X$ 的两个不同实现。这一差值总可由一般函数 $P(y\mid\doxp(x))$ 算出，而后者对 $X$ 的每个水平 $x$ 都有定义，对干预效果提供了更精细的刻画。

\subsection{作为变量的干预}

另一种（有时更有吸引力的）干预说明，把造成干预的力量视为系统内部的一个变量（Pearl 1993b）。其做法是把函数 $f_i$ 本身表示为一个变量 $F_i$ 的取值，并把 \textup{(3.2)} 写成
\begin{equation}
x_i = I(\pa_i, f_i, u_i), \tag{3.7}
\end{equation}
其中 $I$ 是一个三元函数，且只要 $b=f_i$ 就有 $I(a,b,c)=f_i(a,c)$。这相当于把干预概念化为一种外力 $F_i$，它改变 $X_i$ 与其父节点之间的函数 $f_i$。在图形上，我们可以把 $F_i$ 表示为一个新增的 $X_i$ 父节点，而此类干预的效果便可用标准的条件化来分析——即让我们的概率以“变量 $F_i$ 取得取值 $f_i$”这一事件为条件。

原子干预 $\doxp(X_i=x'_i)$ 的效果，通过在 $G$ 中添加一条连线 $F_i\to X_i$ 来编码（见图 3.2），其中 $F_i$ 是一个取值于 $\{\doxp(x'_i),\text{idle}\}$ 的新变量，$x'_i$ 取遍 $X_i$ 的论域，而“idle”表示无干预。于是，$X_i$ 在增广网络中的新父节点集为 $\PA'_i=\PA_i\cup\{F_i\}$，且它与 $X_i$ 通过下述条件概率相联系：
\begin{equation}
P(x_i\mid\pa'_i) =
\begin{cases}
P(x_i\mid\pa_i) & \text{若 } F_i=\text{idle},\\
0 & \text{若 } F_i=\doxp(x'_i)\ \text{且}\ x_i\ne x'_i,\\
1 & \text{若 } F_i=\doxp(x'_i)\ \text{且}\ x_i=x'_i.
\end{cases} \tag{3.8}
\end{equation}
干预 $\doxp(x'_i)$ 的效果，是把原概率函数 $P(x_1,\dots,x_n)$ 变换为新概率函数 $P(x_1,\dots,x_n\mid\hat{x}'_i)$：
\begin{equation}
P(x_1,\dots,x_n\mid\hat{x}'_i) = P'(x_1,\dots,x_n\mid F_i=\doxp(x'_i)), \tag{3.9}
\end{equation}
其中 $P'$ 是由增广网络 $G'=G\cup\{F_i\to X_i\}$ 与 \textup{(3.8)} 所指定的分布，$F_i$ 上带有任意的先验分布。一般地，通过给 $G$ 中每个节点添加一条假设性的干预连线 $F_i\to X_i$，我们可以构造一个增广概率函数 $P'(x_1,\dots,x_n;F_1,\dots,F_n)$，其中含有关于更丰富类型干预的信息。多重干预将由让 $P'$ 以 $F_i$ 的某个子集为条件来表示（这些 $F_i$ 在各自的 $\doxp(x'_i)$ 论域中取值），而干预前的概率函数 $P$ 则可视为把 $P'$ 中每个 $F_i$ 都条件化于“idle”值所诱导的后验分布。

\begin{figure}[htbp]
\centering
\begin{tikzpicture}[causal,scale=0.95,every node/.style={transform shape}]
  \node[vtx] (X) at (-2.6,-0.4) {$X_i$};
  \node[vtx] (P1) at (-3.6,0.7) {};
  \node[vtx] (P2) at (-1.6,0.7) {};
  \node[vtx] (U) at (-2.6,1.7) {$U_i$};
  \node[vtx] (F) at (-0.4,0.6) {$F_i$};
  \draw[ed] (P1)--(X); \draw[ed] (P2)--(X);
  \draw[ed] (U)--(P1); \draw[ed] (U)--(P2);
  \draw[ed,dashed] (F)--(X);
  \node[font=\small,below=0.1mm] at (-2.6,-1.1) {原网络 $G$};
\end{tikzpicture}
\qquad$\Longrightarrow$\qquad
\begin{tikzpicture}[causal,scale=0.95,every node/.style={transform shape}]
  \node[vtx] (X) at (-2.6,-0.4) {$X_i$};
  \node[vtx] (P1) at (-3.6,0.7) {};
  \node[vtx] (P2) at (-1.6,0.7) {};
  \node[vtx] (U) at (-2.6,1.7) {$U_i$};
  \node[vtx] (F) at (-0.4,0.6) {$F_i$};
  \draw[ed] (P1)--(X); \draw[ed] (P2)--(X);
  \draw[ed] (U)--(P1); \draw[ed] (U)--(P2);
  \draw[ed,very thick] (F)--(X);
  \node[font=\small,below=0.1mm] at (-2.6,-1.1) {增广网络 $G'=G\cup\{F_i\to X_i\}$};
\end{tikzpicture}
\caption{用增广网络 $G'=G\cup\{F_i\to X_i\}$ 表示外部干预 $F_i$。}
\label{fig:3.2}
\end{figure}

增广网络表示法的一个优点，是它适用于函数关系 $f_i$ 的任何改变，而不仅限于把 $f_i$ 替换为常数。它还清楚地展示了 $f_i$ 的自发变化（不经外部控制的中介）所造成的推演后果。例如，图 3.2 预测：只有 $X_i$ 的后代会受 $f_i$ 变化的影响，因而对 $X_i$ 非后代的任一集合 $Z$，边缘概率 $P(z)$ 保持不变。同样，图 3.2 规定：对 $X_i$ 后代的任一集合 $Y$，只要 $X_i$ 把 $F_i$ 与 $Y$ d-分离，条件概率 $P(y\mid x_i)$ 对 $f_i$ 的变化保持不变。Kevin Hoover (1990, 2001) 正是利用这类不变特征，通过观察支配经济变量（如就业与货币供给）之过程发生突变（如税制改革、劳资纠纷）所引起的改变，来判定经济变量之间因果影响的方向。确实，每当我们获得可靠信息（例如来自历史或制度知识）表明某个约束给定家族 $(X_i,\PA_i)$ 的特定机制 $f_i$ 发生了突然的局部变化时，就可以利用围绕这些变量观察到的边缘概率与条件概率的变化，来判定 $X_i$ 是否确实是该家族的子节点（或因变量），从而确定该领域中因果影响的结构（Tian and Pearl 2001a）。在此类变化下保持不变的那些统计特征，以及支撑这种不变性的因果假设，都展示在增广网络 $G'$ 中。

\subsection{计算干预的效果}

无论我们把干预表示为一个既有模型的修改（定义 3.2.1），还是表示为增广模型中的条件化（式 \textup{(3.9)}），结果都是干预前分布与干预后分布之间一个良定义的变换。在原子干预 $\doxp(X_i=x'_i)$ 的情形，这一变换可用一个简单的截断分解公式表达，它由 \textup{(3.2)} 与定义 3.2.1 直接得出\footnote{式 \textup{(3.10)} 也可由 Robins (1986, p.~1423) 的 G-计算式（另见 \S 3.6.4）与 Spirtes 等 (1993) 的操纵定理得到（据该文献，此式“由 Fienberg 于 1991 年的一次研讨会上独立猜想”）。关于 \textup{(3.10)} 与 \textup{(3.11)} 所定义之变换的更多性质，见 Goldszmidt and Pearl (1992) 与 Pearl (1993b)。}：
\begin{equation}
P(x_1,\dots,x_n\mid\hat{x}'_i) =
\begin{cases}
\displaystyle\prod_{j\ne i} P(x_j\mid\pa_j) & \text{若 } x_i=x'_i,\\[2mm]
0 & \text{若 } x_i\ne x'_i.
\end{cases} \tag{3.10}
\end{equation}
式 \textup{(3.10)} 反映了从 \textup{(3.5)} 的乘积中移除因子 $P(x_i\mid\pa_i)$，因为 $\pa_i$ 不再影响 $X_i$。例如，干预 $\doxp(X=x')$ 将把 \textup{(3.6)} 所给的干预前分布变换为乘积
\[
P(z_0,z_1,b,z_2,z_3,y\mid\hat{x}') = P(z_0)P(z_1\mid z_0)P(b\mid z_0)P(z_2\mid x',z_1)P(z_3\mid z_2,b)P(y\mid x',z_2,z_3).
\]
在图形上，移除项 $P(x_i\mid\pa_i)$ 等价于移除 $\PA_i$ 与 $X_i$ 之间的连线，同时保持网络其余部分不变。显然，\textup{(3.10)} 所定义的变换满足定义 1.3.1 的条件以及 \textup{(1.38)}—\textup{(1.39)} 的性质。

将 \textup{(3.10)} 乘以并除以 $P(x'_i\mid\pa_i)$，它与干预前分布的关系便更为透明：
\begin{equation}
P(x_1,\dots,x_n\mid\hat{x}'_i) =
\begin{cases}
\displaystyle P(x_1,\dots,x_n)\,\frac{P(x_1,\dots,x_n)}{P(x'_i\mid\pa_i)} & \text{若 } x_i=x'_i,\\[2mm]
0 & \text{若 } x_i\ne x'_i.
\end{cases} \tag{3.11}
\end{equation}
若我们把联合分布看作把“质量”分配给一族抽象点 $(x_1,\dots,x_n)$（每点代表世界的一种可能状态），则 \textup{(3.11)} 所描述的变换揭示了干预 $\doxp(X_i=x'_i)$ 所导致的质量分布变化的一些有趣性质（Goldszmidt and Pearl 1992）。可以看到，每个点 $(x_1,\dots,x_n)$ 的质量都按一个因子增大，该因子等于对应于该点的条件概率 $P(x'_i\mid\pa_i)$ 的倒数。条件概率低的点将大幅提升其质量值；而那些其 $\pa_i$ 取值“预期”了 $X_i=x'_i$ 之自然（非干预）实现的点（即 $P(x'_i\mid\pa_i)\approx1$）则质量不变。在标准的贝叶斯条件化中，每个被排除的点（$x_i\ne x'_i$）都通过一个归一化常数把质量转移给整个被保留点集。然而 \textup{(3.11)} 描述的是一种不同的变换：每个被排除的点（$x_i\ne x'_i$）把质量转移给一组精选的点，这些点与它共享相同的 $\pa_i$ 取值。这可以从以下事实看出：分配给每个层 $\pa_i$ 的总质量、以及每个层内各点的相对质量，都保持不变：
\[
P(\pa_i\mid\doxp(x'_i))=P(\pa_i),\qquad
\frac{P(s_i,\pa_i,x'_i\mid\doxp(x'_i))}{P(s'_i,\pa_i,x'_i\mid\doxp(x'_i))}=\frac{P(s_i,\pa_i,x'_i)}{P(s'_i,\pa_i,x'_i)}.
\]
这里 $S_i$ 表示除 $\{\PA_i\cup X_i\}$ 之外的一切变量。这组精选的受质点，可以视为因共享同一“历史”（由 $\pa_i$ 概括）而与所排除之点“最接近”的点（见 \S 4.1.3 与 \S 7.4.3）。

\textup{(3.11)} 的另一种有趣形式，是把除以 $P(x'_i\mid\pa_i)$ 解释为对 $(X_i=x'_i)$ 与 $\pa_i$ 取条件：
\begin{equation}
P(x_1,\dots,x_n\mid\hat{x}'_i) =
\begin{cases}
P(x_1,\dots,x_n\mid x'_i,\pa_i)\,P(\pa_i) & \text{若 } x_i=x'_i,\\
0 & \text{若 } x_i\ne x'_i.
\end{cases} \tag{3.12}
\end{equation}
当用此式来计算干预 $\doxp(X_i=x'_i)$ 对与 $(X_i\cup\PA_i)$ 不相交的变量集 $Y$ 的效果时，公式便变得熟悉了。把 \textup{(3.12)} 对除 $Y\cup X_i$ 之外的一切变量求和，得到下述定理。

\begin{theorem}[对直接原因作调整]
设 $\PA_i$ 表示变量 $X_i$ 的直接原因集，$Y$ 为与 $\{X_i\cup\PA_i\}$ 不相交的任一变量集。则干预 $\doxp(X_i=x'_i)$ 对 $Y$ 的效果由
\begin{equation}
P(y\mid\hat{x}'_i) = \sum_{\pa_i} P(y\mid x'_i,\pa_i)\,P(\pa_i) \tag{3.13}
\end{equation}
给出，其中 $P(y\mid x'_i,\pa_i)$ 与 $P(\pa_i)$ 表示干预前的概率。
\end{theorem}

式 \textup{(3.13)} 要求先把 $P(y\mid\hat{x}'_i)$ 对 $X_i$ 的父节点取条件，再以 $\PA_i=\pa_i$ 的先验概率为权重对结果取平均。这一取条件再平均的操作称为“对 $\PA_i$ 作调整”。

许多哲学家曾提出这一调整的各种变体，作为因果性与因果效应的概率定义（见 \S 7.5）。例如，Good (1961) 主张对原因发生“直前一刻的宇宙状态”取条件；Suppes (1970) 主张对直到原因发生为止的整个过去取条件；Skyrms (1980, p.~133) 主张对“在决策时刻处于我们影响之外、且与行动结果有因果关联的那些因素的最大限度具体规定”取条件。这些主张中取条件的目的，当然是要消除原因（本例中 $X_i=x'_i$）与结果（$Y=y$）之间的虚假相关；显然，父节点集 $\PA_i$ 能以极大的经济性达成这一目的。在本书所走的\emph{结构论}说明中，因果效应以根本不同的方式定义。条件化算子并非作为旨在压制虚假相关的补救性“调整”而被引入 \textup{(3.13)}；相反，它是从 \textup{(3.10)} 所代表的更深原则——保留干预前分布所能提供的一切不变信息——中形式地浮现出来的。

\textup{(3.10)} 的变换可以容易地推广到同时操纵若干变量的更精细干预。例如，考虑复合干预 $\doxp(S=s)$（$S$ 为变量子集），则（呼应 \textup{(1.37)}）应从 \textup{(3.5)} 的乘积中删去与 $S$ 中变量相对应的全部因子 $P(x_i\mid\pa_i)$，得到更一般的截断分解：
\begin{equation}
P(x_1,\dots,x_n\mid\hat{s}) =
\begin{cases}
\displaystyle\prod_{i\,:\,X_i\notin S} P(x_i\mid\pa_i) & \text{对与 } s \text{ 相容的 } x_1,\dots,x_n,\\[2mm]
0 & \text{其他情形}.
\end{cases} \tag{3.14}
\end{equation}
同样，我们不必局限于把变量设为常数的简单干预。相反，可以考虑对因果模型作更一般的修改，即替换某些机制。例如，若把决定 $X_i$ 取值的机制替换为另一个方程（该方程也许涉及新的变量集 $\PA_i^*$），则所得分布可通过把因子 $P(x_i\mid\pa_i)$ 替换为新方程所诱导的条件概率 $P^*(x_i\mid\pa_i^*)$ 而得到。修改后的联合分布于是为
\[
P^*(x_1,\dots,x_n)=P(x_1,\dots,x_n)\,\frac{P^*(x_i\mid\pa_i^*)}{P(x_i\mid\pa_i)}.
\]

\subsection*{例：动态过程控制}

为说明这些运算，让我们考虑一个过程控制的例子；在健康管理、经济政策制定、产品营销或机器人运动规划等领域的类似应用，应当可以直截了当地照搬。设变量 $Z_k$ 表示生产过程在时刻 $t_k$ 的状态，$X_k$ 表示（在时刻 $t_k$）用于控制该过程的一组变量（见图 3.3）。例如，$Z_k$ 可以表示厂内各处温度与压力之类的测量，$X_k$ 可以表示各种化学品在关键管道中允许流入的速率。假定数据是在过程由策略 $S$ 控制时采集的，其中每个 $X_k$ 由以下两步确定：(i) 监测此前三个变量（$X_{k-1}$、$Z_k$、$Z_{k-1}$）；(ii) 以概率 $P(x_k\mid x_{k-1},z_k,z_{k-1})$ 选择 $X_k=x_k$。

策略 $S$ 的表现被监测并概括为联合概率函数 $P(y,z_1,z_2,\dots,z_n,x_1,x_2,\dots,x_n)$，其中 $Y$ 是结果变量（如最终产品的质量）。

\begin{figure}[htbp]
\centering
\begin{tikzpicture}[causal,scale=1.0,every node/.style={transform shape}]
  \node[vtx] (Z1) at (0,0) {$Z_1$};
  \node[vtx] (X1) at (1.3,1.0) {$X_1$};
  \node[vtx] (Z2) at (2.6,0) {$Z_2$};
  \node[vtx] (X2) at (3.9,1.0) {$X_2$};
  \node[vtx] (dots) at (5.2,0) {$\cdots$};
  \node[vtx] (Zn) at (6.5,0) {$Z_n$};
  \node[vtx] (Xn) at (7.8,1.0) {$X_n$};
  \node[vtx] (Y)  at (9.2,0) {$Y$};
  \draw[ed] (Z1)--(X1); \draw[ed] (X1)--(Z2); \draw[ed] (Z2)--(X2);
  \draw[ed] (X2)--(dots); \draw[ed] (dots)--(Zn); \draw[ed] (Zn)--(Xn);
  \draw[ed] (Xn)--(Y);
  \draw[ed] (Z1)--(Z2); \draw[ed] (X1)--(X2);
  \draw[ed] (X2)--(Zn); \draw[ed] (Zn)--(Y);
  \draw[ed] (X1) to[bend left=18] (Z2);
  \draw[ed] (X2) to[bend left=18] (Zn);
\end{tikzpicture}
\caption{动态因果图，示意序贯过程的控制变量 $X_1,\dots,X_n$、状态变量 $Z_1,\dots,Z_n$ 与结果变量 $Y$ 之间的典型依赖关系。}
\label{fig:3.3}
\end{figure}

最后（为简单起见）假定过程状态 $Z_k$ 只依赖于前一状态 $Z_{k-1}$ 与前一控制 $X_{k-1}$。我们希望评估用新策略 $S^*$ 替换 $S$ 的优劣，新策略依据新的条件概率 $P^*(x_k\mid x_{k-1},z_k,z_{k-1})$ 来选择 $X_k$。

依据此前的分析（式 \textup{(3.14)}），新策略 $S^*$ 的表现 $P^*(y)$ 将由下述分布支配：
\begin{equation}
\begin{aligned}
P^*(y) &= \sum_{z_1,\dots,z_n,x_1,\dots,x_n} P^*(y,z_1,z_2,\dots,z_n,x_1,x_2,\dots,x_n)\\
&= P^*(y\mid z_1,z_2,\dots,z_n,x_1,x_2,\dots,x_n)\prod_{k=1}^{n}P^*(z_k\mid z_{k-1},x_{k-1})\prod_{k=1}^{n}P^*(x_k\mid x_{k-1},z_k,z_{k-1}).
\end{aligned} \tag{3.15}
\end{equation}
由于前两项保持不变、第三项已知，我们有
\begin{equation}
\begin{aligned}
P^*(y) = {}& \sum_{z_1,\dots,z_n,x_1,\dots,x_n} P(y\mid z_1,z_2,\dots,z_n,x_1,x_2,\dots,x_n)\\
&\times \prod_{k=1}^{n}P(z_k\mid z_{k-1},x_{k-1})\prod_{k=1}^{n}P^*(x_k\mid x_{k-1},z_k,z_{k-1}).
\end{aligned} \tag{3.16}
\end{equation}
在 $S^*$ 为确定性且时不变的特别情形，$X_k$ 成为 $X_{k-1}$、$Z_k$、$Z_{k-1}$ 的函数：
\[
x_k = g(x_{k-1},z_k,z_{k-1}).
\]
于是对 $x_1,\dots,x_n$ 的求和可以完成，得到
\begin{equation}
P^*(y) = \sum_{z_1,\dots,z_n} P(y\mid z_1,z_2,\dots,z_n,g_1,g_2,\dots,g_n)\prod_{k=1}^{n}P(z_k\mid z_{k-1},g_{k-1}), \tag{3.17}
\end{equation}
其中 $g_k$ 递归地定义为 $g_1=g(z_1)$ 及 $g_k=g(g_{k-1},z_k,z_{k-1})$。

在 $S^*$ 由若干基本行动 $\doxp(X_k=x_k)$ 构成的特别情形，函数 $g$ 退化为常数 $x_k$，我们得到
\begin{equation}
P^*(y) = P(y\mid\hat{x}_1,\hat{x}_2,\dots,\hat{x}_n) = \sum_{z_1,\dots,z_n} P(y\mid z_1,z_2,\dots,z_n,x_1,x_2,\dots,x_n)\prod_{k}P(z_k\mid z_{k-1},x_{k-1}), \tag{3.18}
\end{equation}
它也可由 \textup{(3.14)} 得到。

这个例子所说明的规划问题，是马尔可夫决策过程（MDP）的典型情形（Howard 1960; Dean and Wellman 1991; Bertsekas and Tsitsiklis 1996），其分析目标是：在给定当前状态 $Z_k$ 与既往行动的条件下，寻找最佳的下一行动 $\doxp(X_k=x_k)$。在 MDP 中，通常给出转移函数 $P(z_{k+1}\mid z_k,x_k)$ 与待最小化的代价函数。在我们刚才分析的问题中，这两个函数都未给出；相反，它们必须从过去（大概是次优的）策略下采集的数据中学习。所幸，由于模型中一切变量都被测量，两个函数都是可识别的，并可直接由相应的条件概率估计如下：
\[
\begin{aligned}
P(z_{k+1}\mid z_k,\hat{x}_k) &= P(z_{k+1}\mid z_k,x_k),\\
P(y\mid z_1,z_2,\dots,z_n,\hat{x}_1,\hat{x}_2,\dots,\hat{x}_n) &= P(y\mid z_1,z_2,\dots,z_n,x_1,x_2,\dots,x_n).
\end{aligned}
\]
第四章（\S 4.4）将处理部分可观测马尔可夫决策过程（POMDP），其中某些状态 $Z_k$ 未被观测；在那些问题中学习转移函数与代价函数，将需要更为精细的识别方法。

值得指出，在这个例子中，为预测新策略的效果，必须先测量那些受某些控制变量（$X_{k-1}$）影响的变量（$Z_k$）。这类测量在实验设计的经典文献中通常被回避（Cox 1958, p.~48），因为它们位于处理与结果之间的因果通路上，容易混淆所需的效应估计。然而我们的分析表明，经过恰当处理，这类测量对于预测动态控制方案的效果可能是不可或缺的。这在与半马尔可夫模型（即含未测量变量的 DAG）有关的情形中尤为真确，后者将在 \S 3.3.2 中分析。

\begin{summarybox}
本节分析的直接含义是：给定一个因果图，其中被干预变量的全部直接原因（即父节点）都可观测，则可由干预前分布推断干预后分布；因此，在此类假设下，我们可以用 \textup{(3.14)} 的截断分解公式，从被动（即非实验）观测中估计干预的效果。然而更具挑战性的问题，是在图 3.1 那样的情形中导出因果效应——那里 $\PA_i$ 的某些成员不可观测，从而妨碍对 $P(x'_i\mid\pa_i)$ 的估计。\S 3.3 与 \S 3.4 将给出简单的图形检验，用以判定 $P(x_j\mid\hat{x}_i)$ 在此类模型中何时可估计。但首先，我们需要更形式地定义“因果量 $Q$ 可由被动观测估计”意味着什么——这个问题属于“识别”（identification）这一术语的范畴。
\end{summarybox}

\subsection{因果量的识别}

因果量不同于统计参数：它们相对于因果模型 $M$ 定义，而不是相对于被观测变量集 $V$ 上的联合分布 $P_M(\cdot)$ 定义。既然非实验数据只提供关于 $P_M(v)$ 的信息，而若干模型又能生成同一分布，就存在这样的危险：所关心的量无法从数据中无歧义地识别出来——即便取无穷多样本亦然。可识别性保证：$M$ 所传达的附加假设（如因果图或结构方程中的零系数）将提供缺失的信息，而不必把 $M$ 完整详尽地阐明。

\begin{definition}[可识别性]
设 $Q(M)$ 是模型 $M$ 的任一可计算量。若对模型类 $\mathcal{M}$ 中任意一对模型 $M_1,M_2$，只要 $P_{M_1}(v)=P_{M_2}(v)$ 就有 $Q(M_1)=Q(M_2)$，则称 $Q$ 在 $\mathcal{M}$ 中\emph{可识别}。若我们的观测有限、只允许估计 $P_M(v)$ 的一个部分特征集 $F_M$，则当 $F_{M_1}=F_{M_2}$ 时总有 $Q(M_1)=Q(M_2)$，便称 $Q$ 可由 $F_M$ \emph{识别}。
\end{definition}

可识别性对于把统计数据（由 $P(v)$ 概括）与关于 $\{f_i\}$ 的不完整因果知识结合起来至关重要，因为它使我们能够从 $P(v)$ 的大样本中一致地估计量 $Q$，而不必具体指明 $M$ 的细节；$\mathcal{M}$ 的一般特征便已足够。就我们的分析而言，所关心的量 $Q$ 是因果效应 $P_M(y\mid\hat{x})$；它当然可由给定的模型 $M$ 算出（用定义 3.2.1），但我们常常需要从 $M$ 的不完整指定——即与 $M$ 相联系的图 $G$ 所刻画的定性特征——来计算它。因此，我们考虑具有下列共同特征的模型类 $\mathcal{M}$：
\begin{enumerate}[label=(\roman*),leftmargin=2.2em]
\item 它们共享相同的子—父家族（即相同的因果图 $G$）；且
\item 它们在可观测变量上诱导出正分布（即 $P(v)>0$）。
\end{enumerate}
相对于这样的类，我们有下述定义。

\begin{definition}[因果效应的可识别性]
若量 $P(y\mid\hat{x})$ 可由被观测变量的任一正概率唯一算出——即对任意两个模型 $M_1,M_2$，只要 $P_{M_1}(v)=P_{M_2}(v)>0$ 且 $G(M_1)=G(M_2)=G$，就有 $P_{M_1}(y\mid\hat{x})=P_{M_2}(y\mid\hat{x})$——则称 $X$ 对 $Y$ 的因果效应可由图 $G$ \emph{识别}。
\end{definition}

$P(y\mid\hat{x})$ 的可识别性保证：可以由两类信息推断行动 $\doxp(X=x)$ 对 $Y$ 的效果：(i) 被动观测，由概率函数 $P(v)$ 概括；(ii) 因果图 $G$，它（定性地）指明哪些变量构成本领域中稳定的机制，或者等价地，指明哪些变量参与决定领域中的每个变量。

把可识别性限制于正分布，可确保条件 $X=x'$ 在数据中以恰当语境得到表示，从而避免 \textup{(3.11)} 中出现零分母。若在行动所施加的语境中 $X$ 从未取到 $x'$，就不可能从这样的数据推断 $\doxp(X=x')$ 的效果。向某些非正分布的推广是可行的，此处不予处理。注意，为证明不可识别性，只需给出两组结构方程，它们在可观测变量上诱导出相同的分布，却具有不同的因果效应。

借助可识别性概念，我们现在可以把 \S 3.2.3 的结果概括为下述定理。

\begin{theorem}
给定任一马尔可夫模型的因果图 $G$，其中变量子集 $V$ 被测量，则只要 $\{X\cup Y\cup\PA_X\}\subseteq V$——即只要 $X$、$Y$ 以及 $X$ 中变量的全部父节点都被测量——因果效应 $P(y\mid\hat{x})$ 便是可识别的。此时 $P(y\mid\hat{x})$ 的表达式由对 $\PA_x$ 作调整而得，如 \textup{(3.13)} 所示。
\end{theorem}

当假定一切变量都被观测时，定理 3.2.5 有一个特别情形。

\begin{corollary}
给定任一马尔可夫模型的因果图 $G$，其中一切变量都被测量，则对任意两个变量子集 $X$ 与 $Y$，因果效应 $P(y\mid\hat{x})$ 都是可识别的，并可由 \textup{(3.14)} 的截断分解得到。
\end{corollary}

现在我们转向半马尔可夫模型中的识别问题。
